% Options for packages loaded elsewhere
\PassOptionsToPackage{unicode}{hyperref}
\PassOptionsToPackage{hyphens}{url}
%
\documentclass[
  11pt,
]{article}
\usepackage{lmodern}
\usepackage{amssymb,amsmath}
\usepackage{ifxetex,ifluatex}
\ifnum 0\ifxetex 1\fi\ifluatex 1\fi=0 % if pdftex
  \usepackage[T1]{fontenc}
  \usepackage[utf8]{inputenc}
  \usepackage{textcomp} % provide euro and other symbols
\else % if luatex or xetex
  \usepackage{unicode-math}
  \defaultfontfeatures{Scale=MatchLowercase}
  \defaultfontfeatures[\rmfamily]{Ligatures=TeX,Scale=1}
\fi
% Use upquote if available, for straight quotes in verbatim environments
\IfFileExists{upquote.sty}{\usepackage{upquote}}{}
\IfFileExists{microtype.sty}{% use microtype if available
  \usepackage[]{microtype}
  \UseMicrotypeSet[protrusion]{basicmath} % disable protrusion for tt fonts
}{}
\makeatletter
\@ifundefined{KOMAClassName}{% if non-KOMA class
  \IfFileExists{parskip.sty}{%
    \usepackage{parskip}
  }{% else
    \setlength{\parindent}{0pt}
    \setlength{\parskip}{6pt plus 2pt minus 1pt}}
}{% if KOMA class
  \KOMAoptions{parskip=half}}
\makeatother
\usepackage{xcolor}
\IfFileExists{xurl.sty}{\usepackage{xurl}}{} % add URL line breaks if available
\IfFileExists{bookmark.sty}{\usepackage{bookmark}}{\usepackage{hyperref}}
\hypersetup{
  hidelinks,
  pdfcreator={LaTeX via pandoc}}
\urlstyle{same} % disable monospaced font for URLs
\usepackage[margin=1in]{geometry}
\setlength{\emergencystretch}{3em} % prevent overfull lines
\providecommand{\tightlist}{%
  \setlength{\itemsep}{0pt}\setlength{\parskip}{0pt}}
\setcounter{secnumdepth}{-\maxdimen} % remove section numbering
\usepackage{esint}
\setmathfont{Asana Math}
\setlength{\emergencystretch}{3em}

\author{}
\date{}

\begin{document}

\hypertarget{bilateral-monopole-saturation-and-exact-bidisk-realization}{%
\section{Bilateral monopole saturation and exact bidisk
realization}\label{bilateral-monopole-saturation-and-exact-bidisk-realization}}

24 September 2026. A continuation of the actual-form and
compact-realization arguments. The new feature is the removal of
\textbf{every topological no-bubble hypothesis}: saturation in both
orientations eliminates the complete Weyl curvature on a harmonic-radius
blowup. The resulting two global parallel forms give an exact locally
symmetric metric on the original smooth manifold. This is a composition
theorem, not a new construction of a bidisk quotient or a claim of
unrestricted Yau26 realization.

\hypertarget{conventions-and-theorem}{%
\subsection{1. Conventions and theorem}\label{conventions-and-theorem}}

Let \(M\) be a closed connected oriented smooth four-manifold. Write
\(M_+=M\) and \(M_-=-M\). A Spin\(^c\) structure on \(M_\sigma\) is
called an \textbf{actual conformal input for \(g\)} here if its
unperturbed Seiberg--Witten equations are solvable for every metric in
\([g]\). The structure is fixed while the positive conformal weight
varies. It may depend on the input metric. The two structures below are
independent; no common spinor, connection, or equation between the
orientations is postulated.

Let \(a_\sigma\) be their real first Chern classes, and set \[
 p_+=\|\operatorname{proj}_{\mathcal H^+_g}a_+\|_{L^2}^2,
 \qquad
 p_-=\|\operatorname{proj}_{\mathcal H^-_g}a_-\|_{L^2}^2.
 \tag{1}
\] Thus \(p_-\) is positive energy, not its negative square in the
original orientation. All norms on \(W^\pm:\Lambda^\pm\to\Lambda^\pm\)
are Hilbert--Schmidt operator norms, norms on two-forms are
exterior-algebra norms, and Ricci uses the tensor norm. Define \[
 \delta=\|\operatorname{Ric}_g^\sharp+I\|_{\infty,\mathrm{op}},
 \qquad r=\fint_M|\operatorname{Ric}_g+g|^2\le4\delta^2,
 \qquad V=\operatorname{Vol}(g),
 \tag{2}
\] \[
 \mathfrak D_\sigma
 =\frac32\int_M(s^2+12|W^\sigma|^2)-72\pi^2p_\sigma,
 \qquad q_\sigma=\mathfrak D_\sigma/V\ge0.
 \tag{3}
\] Nonnegativity is the classical monopole curvature inequality in the
respective orientation.

\textbf{Theorem A (bilateral realization).} For every \(D,A_0>0\) there
exists \(\varepsilon(D,A_0)>0\) with the following property. Suppose
\(M,g\) has actual conformal inputs in both orientations, with \[
 p_+>0,\quad p_->0,\quad \max\{p_+,p_-\}\ge A_0,
 \qquad \operatorname{diam}(g)\le D,
 \tag{4}
\] and \[
                    \delta+q_++q_-<\varepsilon(D,A_0).
 \tag{5}
\] Then the \textbf{same smooth manifold} admits a smooth metric \(h\)
with \(\operatorname{Ric}_h=-h\) whose universal Riemannian cover is \[
                    \mathbb H^2_{-1}\times\mathbb H^2_{-1}.
 \tag{6}
\] Moreover, \(h\) admits global parallel compatible complex structures
\(J_+,J_-\) inducing the two opposite orientations. In particular it is
negative Kähler--Einstein in both orientations.

There is no hypothesis on \(\pi_1(M)\), \(\pi_2(M)\), an ALE core, the
intersection form, or an a priori complex or symplectic structure. No
separate volume lower bound, Weyl bound, curvature-energy bound, Sobolev
constant, class-square cost, or Ricci-derivative bound is required. The
threshold is uniform over varying manifolds and actual class choices.

The conclusion is not that \(g\) is already Einstein, nor that \(M\) is
a globally marked direct product of two surfaces. A compact quotient of
the bidisk need not be a direct product. The actual inputs and smallness
of both deficits remain real restrictions. In particular, this theorem
is not the earlier one-sided theorem with all conditions simply removed.

If either actual integral Chern class has positive square in its own
orientation, its projection has \(p_\sigma\ge1\). Then one may set
\(A_0=1\). No negative-square bound on the other class is introduced.

\hypertarget{the-inherited-gauge-estimates-and-the-volume-regime}{%
\subsection{2. The inherited gauge estimates and the volume
regime}\label{the-inherited-gauge-estimates-and-the-volume-regime}}

For each orientation put \(\kappa_\sigma=r+q_\sigma\). The
\href{../mathematics/current/line1_composed_application.md}{actual-form
theorem}, Sections 2--3, gives, when \(\kappa_\sigma\le1/2880\), a
smooth global form \(\omega_\sigma\in\Omega^\sigma(M)\) for every
\(\eta>0\) sufficiently small, such that \[
 \|\omega_\sigma\|_\infty\le2\sqrt3,
 \qquad
 \fint(|\omega_\sigma|^2-2)^2\le720\kappa_\sigma+\eta,
 \tag{7}
\] \[
 \fint|\nabla\omega_\sigma|^2
 \le13\kappa_\sigma+4\sqrt{2\kappa_\sigma}+\eta.
 \tag{8}
\] These are its zeroth- and first-order conclusions. They have no
negative-square cost. The second-order estimate from that theorem is not
used. In particular, the proof below does not implicitly add the budget
that is needed for a different weak-Bianchi application.

The same classical completed-square estimate, before constructing forms,
gives, with \(w_\sigma=\sqrt6|W^\sigma|\), \[
 \fint(w_\sigma-2)^2\le q_\sigma+2r\le2\kappa_\sigma.
 \tag{9}
\] Indeed \(\int(s+2w_\sigma)^2\le2\mathfrak D_\sigma\), whereas
\(\fint(s+4)^2\le4r\), and \(w_\sigma-2=[(s+2w_\sigma)-(s+4)]/2\).

Consider a sequence satisfying (4) and \(\delta_j+q_{+,j}+q_{-,j}\to0\).
For late terms, \(\delta_j\le1/2\), so
\(-3g_j\le\operatorname{Ric}_{g_j}\le3g_j\). Use the scalar monopole
inequality in whichever orientation has \(p_\sigma\ge A_0\): \[
 32\pi^2 A_0\le\int s_-^2
 \le16(1+\delta_j)^2V_j,
 \qquad V_j\ge v_0:=8\pi^2A_0/9.
 \tag{10}
\] The two choices of orientation give the same scalar curvature and
volume. Bishop--Gromov comparison and the diameter bound consequently
give \[
 V_j\le V_0(D),\qquad
 c(D,A_0)\rho^4\le\operatorname{Vol}_{g_j}B_\rho(x)
 \le C\rho^4\quad(0<\rho\le1).
 \tag{11}
\] For example \(V_0(D)=V_{-1}(D+1)\) suffices; the lower estimate
follows by comparing with the entire manifold inside a \((D+1)\)-ball.

The elementary inequality \(w^2/6\le (w-2)^2/3+4/3\) turns (9) into, for
every measurable set \(E\subset M_j\), \[
 \int_E|W^\sigma_{g_j}|^2
 \le\frac43\operatorname{Vol}_{g_j}(E)
       +\frac23\kappa_{\sigma,j}V_j.
 \tag{12}
\] This holds in \textbf{both} orientations. That is the information
absent from the one-sided realization theorem.

\hypertarget{topology-free-exclusion-of-harmonic-radius-collapse}{%
\subsection{3. Topology-free exclusion of harmonic-radius
collapse}\label{topology-free-exclusion-of-harmonic-radius-collapse}}

\textbf{Lemma 1 (topology-free critical compactness).} Fix \(K,D,v>0\)
and a modulus \(\eta(\rho)\to0\) as \(\rho\downarrow0\). A family of
closed connected smooth four-manifolds satisfying \[
 |\operatorname{Ric}|\le K,\qquad\operatorname{diam}\le D,
 \qquad V\ge v,\qquad
 \sup_x\int_{B_\rho(x)}(|W^+|^2+|W^-|^2)\le\eta(\rho)
 \tag{13a}
\] has a uniform positive \(C^{1,\alpha}\) harmonic radius for each
fixed \(\alpha<1\). It is precompact modulo diffeomorphisms in
\(C^{1,\beta}\), \(\beta<1\), and weak finite \(W^{2,p}\), and has
finitely many smooth types. A sequence in this family with
\(\operatorname{Ric}+g\to0\) uniformly has a smooth negative Einstein
limit on its eventual smooth type and strong finite \(W^{2,p}\)
convergence. No topological, ALE, or separately supplied
curvature-energy hypothesis enters this lemma.

For the sequence in Section 2, (11)--(12) give a common modulus in
(13a). Given a target error, choose a tail for which the
\(\kappa_{\sigma,j}\) terms are small, then choose \(\rho\) small for
its volume terms; the remaining finitely many smooth metrics also have
vanishing small-ball energies. The proof below gives both the lemma and
a direct proof for the sequence, without naming that modulus.

Fix \(0<\alpha<1\) and a standard \(C^{1,\alpha}\) harmonic-radius
definition with fixed metric-comparability and scaled norm bounds.
Suppose the global minimum \(\rho_j\) tends to zero, attained at
\(x_j\). Rescale \[
 \widehat g_j=\rho_j^{-2}g_j,
 \qquad r_h(\widehat g_j,x_j)=1,
 \qquad r_h(\widehat g_j,x)\ge1.
 \tag{13}
\] The rescaled Ricci norm is \(O(\rho_j^2)\), and (11) gives a uniform
Euclidean lower growth estimate on every fixed rescaled ball. The
rescaled manifolds are complete, their diameters tend to infinity, and
the normalized harmonic charts give a pointed complete limit, initially
in \(C^{1,\beta}\), \(\beta<\alpha\).

For clarity, the regularity improvement uses the actual Ricci tensors,
not their derivatives. In a harmonic chart the metric equation is \[
 g^{ab}\partial_a\partial_b g_{ij}
       =-2\operatorname{Ric}_{ij}+Q_{ij}(g,\partial g).
 \tag{14}
\] Its limit is weakly Ricci-flat and hence smooth by elliptic
bootstrapping. Subtract the equation of this smooth limit from (14). The
coefficients converge in \(C^1\), the quadratic lower-order terms
converge, and the Ricci tensors converge uniformly to zero. Interior
elliptic estimates therefore give strong \(W^{2,p}_{\mathrm{loc}}\)
convergence for every finite \(p\). Taking \(p>4/(1-\alpha)\) recovers
strong \(C^{1,\alpha}\) convergence on smaller charts. Curvature
consequently converges strongly in local finite \(L^p\) spaces.
Compatible harmonic charts give the usual embeddings of compact subsets
into late terms. These are the same harmonic-coordinate steps spelled
out in \href{20260923_compact_realization.md}{the compact-realization
proof}, Section 4.

For any fixed \(R\), four-dimensional scale invariance and (12) give \[
 \int_{B_R(x_j,\widehat g_j)}|W^\sigma_{\widehat g_j}|^2
 \le \frac43C(R\rho_j)^4
       +\frac23V_0(D)\kappa_{\sigma,j}\longrightarrow0
       \quad(\sigma=+,-).
 \tag{15}
\] Thus \textbf{both} Weyl halves of the Ricci-flat limit vanish. The
limit is flat; for the general lemma, (15) is replaced by the bound
\(\eta(R\rho_j)\to0\), with exactly the same conclusion. A complete
smooth flat four-manifold with positive Euclidean four-dimensional
volume growth is \(\mathbb R^4\): a nontrivial infinite flat deck group
reduces the volume-growth degree, whereas a finite group of Euclidean
isometries fixes a point and cannot act freely.

On a Euclidean ball of radius larger than two the harmonic-coordinate
defining bounds have strict margins. Strong \(W^{2,p}\) convergence,
followed by solving the harmonic-coordinate Dirichlet problem on a
slightly smaller ball, transfers a chart of radius greater than one to
\(x_j\). This contradicts (13). Hence the original harmonic radii are
bounded below along the sequence.

This argument does not produce a nonflat ALE bubble and then try to
classify or forbid its core. It makes every putative blowup \textbf{flat
before any ALE theorem is invoked}. Neither a global full-curvature
energy estimate nor Jiang--Naber, Bando--Kasue--Nakajima, an ALE
classification, or a hypothesis on the smooth topology is needed here.

The diameter and harmonic-radius bounds now give a finite-atlas compact
limit on a closed smooth manifold. After a subsequence and
orientation-preserving diffeomorphisms \(\phi_j:N\to M_j\), \[
 h_j:=\phi_j^*g_j\longrightarrow h
 \quad\text{in }C^{1,\alpha}\text{ and strong finite }W^{2,p},
 \qquad\operatorname{Ric}_h=-h.
 \tag{16}
\] The same elliptic argument proves smoothness of \(h\); it does not
claim smooth convergence of the original approximate metrics.
Smooth-type preservation is the ordinary harmonic-atlas compactness
conclusion, not a homeomorphism of a possibly singular limit.

\hypertarget{two-global-forms-force-exact-product-geometry}{%
\subsection{4. Two global forms force exact product
geometry}\label{two-global-forms-force-exact-product-geometry}}

Choose \(\eta_j\downarrow0\) and forms from (7)--(8) in both
orientations. Pull them back by the \textbf{same} diffeomorphisms in
(16). Uniform boundedness, the vanishing covariant derivatives in (8),
and \(h_j\to h\) in \(C^1\) make the two sequences bounded in
\(H^1(N,h)\). After a common subsequence they converge strongly in
\(L^2\) to global forms \(\omega_+,\omega_-\). The connection
differences tend uniformly to zero; hence their distributional covariant
derivatives for \(h\) vanish. Elliptic regularity and the amplitude
estimates give \[
 \nabla^h\omega_\pm=0,
 \qquad *_h\omega_\pm=\pm\omega_\pm,
 \qquad|\omega_\pm|_h^2=2.
 \tag{17}
\] In particular, neither form vanishes. There is no eigenline sign
choice and no two-fold cover: the approximating forms were global.

Define \(J_\pm\) by \(h(J_\pm X,Y)=\omega_\pm(X,Y)\). In four dimensions
(17) gives \(J_\pm^2=-I\). Self-dual and anti-self-dual skew
endomorphisms commute, so \[
 K=-J_+J_-,\qquad K^*=K,\qquad K^2=I,\qquad
 \operatorname{tr}K=0,\qquad\nabla K=0.
 \tag{18}
\] The trace follows equally from the orthogonality of the two forms.
Consequently \(TN=E_1\oplus E_2\), the two orthogonal rank-two parallel
eigenbundles of \(K\). Parallel distributions are integrable and their
leaves are totally geodesic. On the complete simply connected universal
cover the de Rham splitting is a global Riemannian product of two
complete simply connected surfaces. Since \(\operatorname{Ric}_h=-h\),
the Ricci tensor of each factor is minus its metric. In dimension two
this means constant Gaussian curvature \(-1\). Both factors are
\(\mathbb H^2_{-1}\), proving (6). Parallelness of \(J_\pm\) also proves
their integrability and the global Kähler assertions directly.

If Theorem A failed for fixed \(D,A_0\), choose counterexamples with
\(\delta_j+q_{+,j}+q_{-,j}<1/j\). Sections 2--4 give (16)--(18). Pulling
\(h,J_+,J_-\) forward to each late \(M_j\) contradicts its being a
counterexample. This proves a fixed threshold uniform across changing
smooth types. It does not assume that an individual small-error metric
can be improved into a whole sequence on its own manifold. \(\square\)

\hypertarget{raw-ricci-energy-formulation}{%
\subsection{5. Raw Ricci-energy
formulation}\label{raw-ricci-energy-formulation}}

Let \[
 T_\pm=2\chi(M)\pm3\tau(M),\quad
 R_2=\int_M|\operatorname{Ric}|^2,\quad
 S_2=\int_Ms^2,\quad P=\int_M|\operatorname{Ric}^{\circ}|^2.
\] Since \(R_2=P+S_2/4\) and
\(4\pi^2T_\sigma=S_2/24+2\int|W^\sigma|^2-P/2\), \[
 \boxed{
 \mathfrak D_\sigma
   =\frac92\bigl[R_2-8\pi^2(2p_\sigma-T_\sigma)\bigr].}
 \tag{19}
\] This is a classical Chern--Weil identity, the exact deficit form of
LeBrun's Ricci-energy estimate. It is not a new inequality.

Suppose independent actual projection budgets are available,
\(p_\sigma\ge A_\sigma>0\), and \(\max(A_+,A_-)\ge A_0\). Define the
nonnegative averaged two-sided energy excess \[
 e_{\rm bi}
 =\frac{R_2-8\pi^2(A_++A_--2\chi)}{V}.
 \tag{20}
\] Adding (19) and using \(T_++T_-=4\chi\) gives the exact decomposition
\[
 \boxed{
 e_{\rm bi}
 =\frac{q_++q_-}{9}
   +\frac{8\pi^2}{V}
      \bigl[(p_+-A_+)+(p_--A_-)\bigr].}
 \tag{21}
\] Every term on the right is nonnegative. Thus no cancellation between
the orientations is hidden by the single energy condition.

\textbf{Corollary B (energy-to-geometry realization).} Under the
actual-input, budget and diameter assumptions just stated, without any
topological restriction, the condition \[
                  \delta+9e_{\rm bi}<\varepsilon(D,A_0)
 \tag{22}
\] implies the full conclusion of Theorem A.

Indeed (21) gives \(q_++q_-\le9e_{\rm bi}\), so (5) follows. This
replaces the two separately displayed functional deficits by a raw
Ricci-energy excess above the average of two genuine gauge bounds. It
does not replace the existence of the two actual inputs by numerical
cohomology data, and it does not claim that (20) is the attained global
Ricci-energy infimum on every supplied topology.

The exact counterpart using the actual projections is \[
 q_++q_-=9\,\frac{R_2-8\pi^2(p_++p_--2\chi)}{V}.
 \tag{23}
\] Neither (22) nor (23) follows from \(\delta\ll1\) alone. Their
purpose is a checkable interface for independent topological or
variational inputs.

\hypertarget{consequences-tests-and-scope}{%
\subsection{6. Consequences, tests, and
scope}\label{consequences-tests-and-scope}}

\hypertarget{topology-is-now-a-conclusion-not-an-assumption}{%
\subsubsection{6.1 Topology is now a conclusion, not an
assumption}\label{topology-is-now-a-conclusion-not-an-assumption}}

Let \(\mathcal B\) denote the smooth manifolds admitting a metric with
universal Riemannian cover (6) and the two global parallel structures of
Theorem A. If \(M\notin\mathcal B\), then every input meeting (4) obeys
\[
                  \delta+q_++q_-\ge\varepsilon(D,A_0).
 \tag{24}
\] Examples of independently checkable exclusions from \(\mathcal B\)
are \(\pi_2(M)\ne0\), torsion in \(\pi_1(M)\), or \(\tau(M)\ne0\).
Indeed (6) is contractible; a finite-order deck isometry of this
Hadamard space has a fixed point and cannot act freely; and for the
product metric both Weyl halves have the same squared norm \(2/3\), so
the signature is zero. The output also has \(\chi(M)>0\).

Thus the theorem does not merely improve a threshold within the old
no-ALE-core class. It tests inputs on \textbf{arbitrary} smooth
four-manifolds and, when they pass, deduces the topology and geometry
that the earlier realization theorem had partly assumed. The new
condition is bilateral saturation; neither theorem subsumes the other on
all its inputs.

Equation (24) is a joint deficit gap, not a pointwise-only gap on every
non-bidisk manifold. No claim follows that such a manifold has no other
Einstein metric.

\hypertarget{actual-equality-models}{%
\subsubsection{6.2 Actual equality
models}\label{actual-equality-models}}

For \(M=\Sigma_g\times\Sigma_h\), \(g,h\ge2\), take the product of
curvature \(-1\) metrics and complex structures \(J_+=(J_1,J_2)\),
\(J_-=(J_1,-J_2)\). The canonical Spin\(^c\) structures in the two
orientations supply actual classes: both are general-type Kähler
surfaces with \(b_\pm>1\), and their canonical Seiberg--Witten
invariants are nonzero. Set \(A=8(g-1)(h-1)=2\chi(M)\). Then \[
 V=2\pi^2A,\quad p_+=p_-=A,
 \quad |W^+|^2=|W^-|^2=2/3,
 \quad \delta=q_+=q_-=e_{\rm bi}=0
 \tag{25}
\] with \(A_+=A_-=A\). This verifies simultaneous saturation in a real
nonempty family. It is a classical equality model, not a newly
constructed family. The theorem's increment is the uniform implication
for arbitrary smooth inputs and changing topologies, without a presumed
product structure.

Unequally scaled hyperbolic products remain cscK in both orientations
and have \(q_+=q_-=0\), while generally \(\delta\ne0\). They test why
simultaneous functional equality must not be identified with the input
metric itself being Einstein. Their underlying manifold already admits
the normalized product metric, consistently with the existence
conclusion.

\hypertarget{diameter-cannot-simply-be-dropped-even-with-both-inputs}{%
\subsubsection{6.3 Diameter cannot simply be dropped, even with both
inputs}\label{diameter-cannot-simply-be-dropped-even-with-both-inputs}}

The fixed nonzero-shear self-gluings of
\href{../mathematics/current/round11_closed_gluing.md}{Round 11},
Section 6, carry symplectic structures in \textbf{both} orientations.
Their canonical classes are actual monopole inputs with
\(p_\pm\ge A=2\chi>0\); this follows from \(b_\pm>1\) and Taubes as
recorded in \href{../mathematics/current/round11_variational.md}{the
variational proof}, Sections 1--2. The actual normalized metrics
\(g_\ell\) satisfy \[
 \tau=0,\quad V_\ell=2\pi^2A,\quad
 \delta_\ell\to0,
 \tag{26}
\] and \href{../mathematics/current/round11_kahler.md}{the global
topology proof} excludes a Kähler structure on every finite cover.

Using (19), \(R_2/V\le4(1+\delta_\ell)^2\), \(T_\pm=A\), and
\(p_\pm\ge A\), one obtains the simultaneous bounds \[
 0\le q_{\pm,\ell}
 \le18\bigl[(1+\delta_\ell)^2-1\bigr]
 =36\delta_\ell+18\delta_\ell^2\longrightarrow0.
 \tag{27}
\] Thus bilateral actual saturation and pointwise almost-Einstein input
do occur on a fixed manifold with no bidisk metric, when diameter is
unrestricted. Theorem A forces \(\operatorname{diam}(g_\ell)\to\infty\).
Equation (27) uses the direct Ricci-energy estimate and is slightly
tighter than the separate scalar/traceless estimates in the one-sided
note; that numerical improvement is not a separate mathematical
achievement.

These are inherited constructions and a new consequence, not a new
counterexample construction. They disprove diameter-free \textbf{bidisk
or KE} realization with these inputs. They do not decide whether those
manifolds carry a general non-Kähler Einstein metric, and therefore do
not furnish Yau26 counterexamples.

\hypertarget{what-is-and-is-not-new-in-the-proof}{%
\subsubsection{6.4 What is and is not new in the
proof}\label{what-is-and-is-not-new-in-the-proof}}

The scalar and curvature-functional bounds, Chern--Weil algebra,
harmonic compactness and de Rham splitting are classical. The actual
global form estimates are inherited first-line results. The new
composition uses the \textbf{second actual orientation} twice: to rule
out all harmonic-radius blowups without topology, and to upgrade the
limiting KE structure to an exact bidisk realization. It thereby
connects two-sided near-optimal Ricci energy to exact geometry on the
same smooth manifold.

The result establishes the Yau26 existence output on its stated inputs;
locally symmetric negative KE is more than the required output. The
remaining limitations are the diameter, actual bilateral class supply,
and simultaneous-saturation inputs. The proof supplies neither a
dimension-only threshold nor an effective numerical constant. No
exhaustive priority claim is made.

\hypertarget{dependencies-and-finite-verification}{%
\subsection{7. Dependencies and finite
verification}\label{dependencies-and-finite-verification}}

\begin{itemize}
\tightlist
\item
  \href{../mathematics/current/line1_composed_application.md}{Inherited
  actual-form theorem}, Sections 2--3: global forms and the explicit
  amplitude/gradient estimates.
\item
  \href{20260923_gauge_saturation_ke.md}{One-sided realization},
  Sections 3--6: the original composition and one-form limit. The
  present argument replaces its topological no-bubbling step, not its
  form supply.
\item
  \href{20260923_compact_realization.md}{Harmonic-coordinate realization
  proof}, Sections 3--5: the metric-equation convergence and
  preservation of smooth type. Its ALE and classification steps are not
  used here.
\item
  \href{https://arxiv.org/html/math/0003068}{LeBrun, \emph{Ricci
  Curvature, Minimal Volumes, and Seiberg--Witten Theory}}, Theorem 2.3
  and Propositions 4.1--4.2: actual-conformal curvature inequality, cscK
  equality, and Ricci-energy identity. The energy reformulation above is
  classical algebra; it is not attributed as a new part of this
  continuation.
\item
  \href{https://doi.org/10.1007/BF01233434}{Anderson, \emph{Convergence
  and rigidity of manifolds under Ricci curvature bounds}}, and
  \href{https://arxiv.org/html/1705.10071}{Zhang--Zhu, \emph{Bounds on
  harmonic radius and limits of manifolds with bounded Bakry--Émery
  Ricci curvature}}, Section 2, especially its construction of larger
  harmonic charts in a Euclidean limit: the harmonic-coordinate
  compactness framework. The Ricci-flat equation and strong
  finite-Sobolev upgrade needed here are written explicitly in Section
  3.
\item
  \href{check_bilateral_realization.py}{Finite check} verifies (19),
  (21), the product normalization, and the commuting opposite-handed
  complex structures. It is not a computer proof of gauge solvability,
  harmonic compactness, or the realization theorem.
\end{itemize}

\end{document}
